\section{从应天到北平：建国仍是一段行程}

洪武元年正月，应天城内的即位仪式使朱元璋能够以皇帝名义颁诏、任官和调兵。这个名义依托的是已经握在手中的江南粮赋、仓储和军政节点；它没有在同一天越过长江、进入北方城镇，或替地方社会决定应当如何响应。南京明故宫的营建遗存让人看见新中枢所需的宫城和工役，而即位诏所能确证的是朝廷发令的起点，不是各地归附已经完成。

同年八月，\eventanchor{EM-1368-DADU}{徐达军进入大都并改称北平}。元顺帝已经北撤，城内易手因而不能被写成一场把整个北方一并交付的胜利。对入城的军队来说，城门、仓储、旧官署和通往河北的道路要逐项接管；对城内居民和旧有吏役来说，改换名号还要落实为谁征粮、谁守城、谁保存原有生计。\textit{太祖实录}记下的是朝廷军事行动的次序，\textit{高丽史}提供北方消息的另一端；二者并不能替城中损失、去留与每一处衙署的交接作证。

\section{草原与海面：两条没有封死的边界（1369--1372）}

\eventanchor{EM-1369-SHANGDU}{洪武二年明军北进上都、应昌一线}，追逐的对象已经不只是大都的城墙。北元皇室和部众在草原上移动，徐达所部必须计算水源、营地、马匹与补给；元顺帝北遁途中去世，又使旧朝的中心继续变动。占据一座都城不能自动占据草原上的政治关系。到\eventanchor{EM-1372-LINGBEI-EXPEDITION}{洪武五年三路北伐受挫}，这种限制变得具体：各路军队难以按原定设想会合，行军距离、气候和给养使京师的意图不能直接变成边地的结果。战役大势可以互校，败因和损失却在记载中不一；较稳妥的结论是，北方经营从此要倚重卫所、边镇与时断时续的交涉，不能只靠一次北伐收束。

另一条边界在东南海面。\eventanchor{EM-1371-HAIJIN}{洪武四年申禁滨海居民私自出海}，把航行、货物和来使纳入官府可辨认的名目。港口并没有因此成为静止的线条：船户要面对季风和船只，商人要寻找货源，沿岸官员也要在禁令、渔盐与治安之间作出判断。其后遣使琉球，恰好说明禁令与官方往来并非彼此取消；朝廷一面压缩未经许可的航行，一面保留能够册封、交涉和贸易的通道。\textit{历代宝案}所见的是琉球一端的收文和遣使，不能把岛内各群体或每一港口的实际停航写成同一种经验。

\section{把机构拆开，把责任写下（1376--1387）}

\eventanchor{EM-1376-PROVINCIAL-COMMISSIONS}{洪武九年废行中书省，分置承宣布政使司、都指挥使司与提刑按察使司}。这项安排不是把一个现代意义的省一分为三，而是要避免民政、军政和监察重新汇在同一位地方长官手中。布政司、都司和按察司各自承接不同文书与人手，遇到征兵、诉讼或军粮时又不得不彼此通报。命令可以确定，地方运转却取决于书吏、驿递、卫所和府县是否把它接续下去；会典中后来归整出的制度轮廓，不能倒过来证明每一处三司立即均衡有效。

\eventanchor{EM-1380-HU-WEIYONG}{洪武十三年胡惟庸案后中书省被废}，把同一种防范推进到南京中枢。六部文书不再经过宰相机构而直接面对皇帝。朝廷由此避免一处职位长期聚集人事与政务，却把原先由中书省承受的协调压力分摊给六部、近臣和皇帝本人。案狱所称的“通倭”“通元”是定罪过程中的官方语言，不能作为无须证明的事实；可确知的是，谁汇总奏报、谁承担迟误和误判的责任，从此更紧密地向宫城收束。

洪武十四年，\eventanchor{EM-1381-LIJIA-HUANGCE}{黄册与里甲编制被推进}。战争、迁徙和流亡已经打乱旧有的居住与役属关系，国家试图以户等把赋役、军户和田土重新接入一条可追查的链。册籍因此不是对稳定乡村的中性摄影，而是催征、协商与规避同时发生的工具。\eventanchor{EM-1387-FISHSCALE-REGISTERS}{洪武二十年催成的鱼鳞图册等田土登记}，又试图把地块、界址和赋额接入核算。黄册偏重人户的责任，鱼鳞图册更细地处理田土；徽州残存册页与契约能显示若干地点如何书写这种关系，却不能拼成全明精确的人口、亩数或实收总额。抄补层次、保存链和地方差异，正是这类材料必须保留的限制。

\section{西南不是北方胜利的附注（1381--1382）}

同在洪武十四年，\eventanchor{EM-1381-YUNNAN-CAMPAIGN}{傅友德、蓝玉、沐英由川黔向云南推进}。军队能否继续前行，要看山路、渡口和粮运能否支撑，也要看沿线的地方势力如何衡量抵抗、避让或交涉。南京的中枢决断至此接受一次远距离的检验：主将要把不同方向的队伍送到曲靖一线，沿途的人则要承担供给、消息和通行的成本。明方奏报可确认出兵与主将，不能把军额、伤亡或各社群的选择化为准确而单一的统计。

\eventanchor{EM-1382-YUNNAN-ANNEXATION}{次年昆明陷落、梁王自尽后，明廷开始在云南布置卫所与建置}。一座城和若干行政节点归入新朝，并不等于山地、城郊与商路都已同样服从。军屯要有粮种和驻军，卫所要有可走的道路，土司关系也需要翻译者、地方首领与官员不断协商。大理、昆明的碑刻和地方志偶尔保留这种接收的痕迹；材料较少之处不能用“全境归附”补满。征服所获得的是延伸国家能力的机会，也留下谁支付长期供给与治理成本的问题。

\section{继承的名分进入战争（1398--1402）}

\eventanchor{EM-1398-HONGWU-DEATH}{洪武三十一年闰五月太祖去世，皇太孙朱允炆继位}。继承日期可以确定，遗诏、宫中部署和诸王即时反应却经过建文、永乐两朝的修史，不能当作透明的记录。成年藩王仍拥有封国、卫所和边地联系，新君需要的不是一张已经获得一致承认的名单，而是一套能让官僚、宗室与军队继续发令的秩序。

\eventanchor{EM-1399-JINGNAN-BEGIN}{建文元年燕王朱棣以“靖难”为名起兵}，使争论落到河北道路、北平卫所、山东州县和运河补给上。“削藩”与“奉天靖难”都是争夺合法性的语言，不能替任何一方预先证明正当性。战争推进时，沿线州县首先感到的是征发、运输和道路损耗；白沟河与东昌的胜负可以互校，军数、火器作用和责任归属却不应照胜者叙事写定。

\eventanchor{EM-1402-NANJING-FALL}{建文四年燕军入南京，朱棣即位}。金川门、宫城火灾和建文帝去向在后出叙事中尤其密集，而入城和改元是能够较稳地定位的节点。取得都城意味着取得印信、仓储、官署和重编前朝记录的条件，却没有立即消除内战留下的人事与地方压力。永乐朝重修建文材料本身，便提醒读者：谁控制中枢，谁也影响后人怎样排列这场继承战争。

\section{北京、文献与海路：经营新中枢（1403--1408）}

\eventanchor{EM-1403-BEIJING-NAMING}{永乐元年北平改称北京}，标志着北方城市被重新安放在王朝的政治与军政计算中。改名不是迁都的完成；宫城、坛庙、仓储、城防、卫戍和南方粮道仍须逐项建造和维持。同年，\eventanchor{EM-1403-YONGLE-DADIAN-COMMISSION}{解缙等受命编纂大型类书}。编纂工程把抄写者、校勘者、藏书与宫廷资源集中到一处，显示新朝也在安排什么文本可以被汇集和分类；现存卷册题记能留下工程的物质痕迹，却不能使“收尽天下书”的宣称变成无遗漏或广泛传播的事实。

\eventanchor{EM-1405-ZHENGHE-FIRST-VOYAGE}{永乐三年，郑和舰队自南京、刘家港出航}。造船、仓储、翻译与沿海补给被同一套官府调度，出海以后却仍要依赖季风、港口准入和各地中介。朝廷的遣使、马欢与费信所记的见闻、以及琉球等地的外交文书，观察的是不同的行动，并不合成一张可以精确量出船数、船型和贸易额的地图。首航的后果不是国家突然控制印度洋，而是让东南的劳力、物资与使团网络承担了更长距离的联系。

\eventanchor{EM-1407-JIAOZHI-PROVINCE}{永乐五年废大虞、设交趾承宣布政使司}，把直接统治的选择推向红河三角洲和山地之间。诏令能显示北京如何命名这个区域，不能证明驻军、官署与税役已经穿透所有地方。越南编年材料与明方记载对战争和建置提供不同观察位置；后来的撤军也提醒人们，设置行政名目不等于解决征粮、补给和地方抵抗。

次年\eventanchor{EM-1408-YONGLE-DADIAN-COMPLETED}{《永乐大典》进呈}，是上述文献工程的一个可见节点。进呈将大量抄写、校勘和贮藏的工作呈到宫廷，但现存残卷与后出的传记只能让我们接近它的编纂和留存，不能替地方读者作证。北京的营建、海外使节与大典的编纂因此不是彼此孤立的“盛事”：它们都要求人力、文书、运输和可长期维持的财政。

\section{北征、运河与西南道路（1410--1417）}

\eventanchor{EM-1410-FIRST-MONGOL-CAMPAIGN}{永乐八年亲征蒙古}，把北京方向的供给问题推到更远的草原。军队沿途需要水源、牧草、车马和补给，敌方部众却并非一支静止等待会战的队伍。亲征和作战可以确认，明方所称战果、敌军人数和部众归属则须与蒙古叙事相互审读。北征增强了中枢对边防的关注，也使北京必须持续取得南方和沿线的资源。

因此\eventanchor{EM-1411-GRAND-CANAL}{永乐九年大规模疏浚会通河}，不能只读作一项水利成就。河工要处理闸坝、河道、役夫与地方征发，航路能否使用还取决于河段维护和粮运组织。\eventanchor{EM-1415-CANAL-TRANSPORT}{到永乐十三年会通河航路投入漕粮北运}，江南粮食才较稳定地转向北京；“较稳定”不是无损耗、无代价。沿线碑刻可以让若干闸坝和工程露出位置，无法据此计算所有民夫负担或每年的实运数额。北方中枢的供给被改善，也更深地把地方财政、劳役与河道风险捆在一起。

\eventanchor{EM-1413-GUIZHOU-PROVINCE}{永乐十一年设置贵州承宣布政使司}，同样显示诏令必须经过山路、驿站、卫所、土司和州县才会成为治理。新机构想重新安排军粮、诉讼和消息的接口，却仍依赖地方首领、翻译者和驻军合作；不能用一个“土司区”抹去不同社群的关系。同年\eventanchor{EM-1413-ZHENGHE-FOURTH-VOYAGE}{郑和第四次下西洋抵达忽鲁谟斯等港口}，航程延长使使团更依赖港口翻译、补给和在地交换。航次与主要目的地可证，具体路线和所见情形仍须区别作者亲历、抄本层次和后出汇编。

\eventanchor{EM-1414-TULA}{永乐十二年北征瓦剌、在土拉河一带交战}，又把草原行动的代价带回京师。明军可将一次交战写为胜利，瓦剌内部损失与部众去向却不能由明方歼获数断定。次年运河通航的意义正正在此：它不是孤立地“便利交通”，而是给京师、北征和营建准备一条必须不断维修的给养线。\eventanchor{EM-1416-BEIJING-CONSTRUCTION}{永乐十四年宫城、坛庙与城防工程持续集中施工}，木材、工匠、粮食和道路于是同时承压；遗址分期可说明工程阶段，不能精确还原每座建筑的完工日。

\section{定都之后：工程、远征与继承（1420--1424）}

\eventanchor{EM-1420-BEIJING-CAPITAL}{永乐十八年正式定北京为京师}，南京仍保有留都地位。这个双都安排不是旧都一夜失去功能，而是中枢常驻、边防决策和北运粮道更紧密地相连。翌年\eventanchor{EM-1421-PALACE-FIRE}{奉天、华盖、谨身三殿失火}，重建需要重新安排木材、工役与仓储；把火灾解释为天意是当时政治修辞，并不能代替对起火原因和损失的证据。

\eventanchor{EM-1422-SECOND-NORTHERN-CAMPAIGN}{永乐二十年再度北征，寻击阿鲁台而未获决定性会战}；\eventanchor{EM-1423-THIRD-NORTHERN-CAMPAIGN}{次年继续北征、搜寻其部众}。出征与行程可从朝廷记录定位，但“无功”或“震慑成功”都是带有政治立场的评语。草原对手的移动迫使军队追逐不断变化的消息，北方道路和补给也因此继续承担成本。最后，\eventanchor{EM-1424-YONGLE-DEATH}{永乐二十二年皇帝北征归途中去世，太子朱高炽继位}。死亡和继位可以确认，秘不发丧及随行谋划的热闹细节多依后出叙事，不能填补为确定的宫廷现场。

\subsection{册籍落到何处：一户人家的位置并不由一张纸决定}

黄册、里甲与鱼鳞图册给朝廷提供了把户、田、役逐项对照的语言，却没有消除它们之间的缝隙。一次迁居会使原籍、现居、军役和田土的写法彼此错开；一份分家契约则可能把同一块土地写成边界、价银、债务和亲属义务，而不是户部需要的赋额。对地方书手、里长和业主而言，造册既是向上报送责任，也是在熟人、田界和追征之间留下可以解释的余地。故而册页上的整齐栏目，不宜被读成乡村已经被完整掌握；它更准确地显示，国家在何种格式中要求人们被看见。

军户的处境尤其说明这一点。卫所需要能够出役的人、屯田需要耕作者，远征或守边又会把家庭带离原来的田地和亲属网络。册上身份使朝廷能够提出调发，不能保证军户始终有足够的劳力、种子或路费完成命令。北方边镇、云南卫所和运河沿线所留下的材料密度很不相同：北京的档案往往看见的是兵额和命令，碑刻、契约、方志和个案才偶尔让某家怎样迁移、逃役或兼营露出轮廓。材料的沉默不是地方没有经历变化，而是许多经历没有以能保存到今天的形式进入中枢。

这一差别也改变我们理解“地方接收”的方式。新朝给某处命名、派出官员或设立卫所，可以证明国家把那里纳入了一条行政设想；实际的征粮、审案、借贷、婚姻和劳作，则由县衙、宗族、市场和地方权威共同塑形。早明国家能力的增长因此既来自可重复的登记和催征，也受制于它必须倚靠的中介。日后朝廷反复更改法令、重核册籍和追问欠额，并非因为原有文本不存在，而是因为文本进入地方以后总会遇到人、地和生计的变化。

\subsection{从南京到交趾：官署越过边界以后}

永乐朝在交趾的选择把这个问题推得更远。红河三角洲的城镇、河道与粮区能够使军队和官署暂时取得立足点，山区的道路、村社与地方首领却不接受一条由北京单向画出的行政线。驻军要粮，文书要译，征税与治安又需要熟悉当地关系的人参与；每一个接口都可能把朝廷的命令改造成协商、拖延、抵制或新的负担。把设省理解为“越南已经成为明朝一部分”，会把直接统治正在面对的这些成本抹掉。

明方实录擅长保存诏令、任命和战报的中枢顺序，越南编年史保留的则是另一种政权如何组织抵抗与记忆战争。两边材料可以帮助确认若干行动的先后，不能消除各自的胜败框架，更不能让读者得知每一处村落怎样经历征发。永乐后期持续送往南方的兵员与物资，正与北京的建设、北征和漕运竞争。多年后撤出交趾不是对一四〇七年建置的简单否定，而是显示一项可以颁布的制度未必能负担持久的道路、粮饷和地方关系。

海上的使团也有类似的行政距离。郑和舰队从南京和刘家港起航时，官府能够安排赏赐、诏书、翻译者和随行人员；船队抵达港口后，却要进入当地统治者、商人、航海者和解释者已经形成的秩序。马欢、费信的书写最有价值之处，不在于替朝廷证明一个无边界的海洋帝国，而在于让港口的翻译、物产、礼仪和交换显露出来。它们也有使团观看异域的格式。读者应当同时看到北京如何发令和港口如何使命令变形，而不是把朝贡名目、私人贸易和实际控制混为同一件事。

\subsection{河道把不同地方放在同一份账上}

会通河的疏浚和通航，使“南粮北运”不再只是朝廷对地图的设想。山东的河段、闸坝与堤防必须容纳水势，沿岸州县要承受河工、夫役和维修，船只、仓储与验收又让米粮在每一个节点可能损耗或停滞。京师得到更可持续的供给，并不表示江南或运河沿线自动从中受益；对当地官员和家庭而言，漕运可以是市场机会，也可以是临时征发、堤防决口和催办任务。不同河段、年度和口径中的起运、报数、实到不能彼此替换。

北京的宫殿、坛庙和城防工程正把这种物流关系固定下来。木材、砖石、工匠与军队的需要使北方都城不是一座只消耗政治象征的城市，而是许多省份和道路共同维持的节点。南京保留留都地位，意味着旧都的官署、仓储和文化资源并未消失；双都格局反而要求文书与物资在更长距离上不断协调。奉天等殿的火灾之所以值得放在迁都之后观察，正在于它迫使朝廷再次决定什么先运、什么先修、由谁承受新的劳役，而非仅仅提供一则灾异故事。

北征亦不能与这些工程分开。皇帝出塞时，草原行动要依赖车马、粮秣与可用消息；对手的游动则使“找到并决战”始终未必能够实现。几次远征所留下的，不只是一串出发与回师的年份，还有北京在战争、营建和漕运之间分配有限资源的日常难题。成祖最后一次北征途中去世，继承人面对的正是一副已经拉得很长的国家骨架：它可以把江南粮食运到北方、把诏令送往山地和海港，却也必须持续支付让这些关系不至断裂的代价。

\subsection{谁能够读到命令，谁能够把命令改写成日常}

新朝还要解决官员从何而来的问题。国子学、学校与取士政策为六部和地方官署准备识字的行政人手，但学校建置本身不能说明各地有同样的师资、经费或入学机会。洪武朝曾在科举、荐举和学校考核之间调整，南北榜案又使籍贯和取士分布进入中枢政治。朝廷能够由榜单和诏令表明它如何处理资格，不能据此断定一地士人、宗族和普通家庭对这种安排有相同的感受。一个人取得资格以后，仍需穿过荐举、任命、驿路和地方官署，才可能成为将政令带到县份的人。

《永乐大典》的编纂与这种官僚问题并行。解缙等人主持的类书工程要搜集、抄写、校勘和分类，依赖的不只是皇帝的诏令，还有能供应纸墨、识字和贮藏的机构。现存卷册使读者看见一项国家工程确曾留下材料，也让人注意到今日所存只是庞大工程的残余。把它当作“知识已经统一”的证据，会忽略文本进入书院、家庭、寺观、市场和边地时仍受抄写成本、读者能力和地方网络限制。朝廷能决定一批文献应当被汇集，却不能替所有人决定何时、怎样读它。

同样的距离存在于法律和案件之间。\textit{大诰}以案件与重刑规定国家希望官民畏惧什么，锦衣卫使皇帝拥有不完全依赖常规部院的侦缉与转达渠道；这首先改变的是谁能够把消息送入宫城、谁会被当作可疑者。它们不能直接换算为犯罪率、基层知法程度或每一次审判的实际过程。征粮、户籍和军役越要求准确，承办这些事务的人越可能在恐惧、自保和地方关系中重新书写事实。明初制度的锐利之处正在这里：它扩大了中枢追问的范围，也扩大了中枢不得不依赖中介的范围。

\subsection{边防并不在长城外才开始}

北方经营的压力并非只在皇帝出塞时显现。卫所、边镇和北平周边的城池需要驻军、马匹、仓粮和可修复的道路；北京成为常驻中枢以后，边情如何呈报、粮饷如何北送、沿线谁负责保管，便进入官员与地方社会的日常。草原部众的称谓和联盟在明方文书中常随战报而变化，不能把“蒙古”写成内部没有差异的单一对手。较能确定的是，明廷要维持一条向北伸出的行动能力，便必须同时维持来自南方的粮运和来自地方的兵役。

这种相互依赖也解释了为什么北征后未获决定性会战仍有制度后果。军队回师并不会自动收回已经消耗的粮秣、劳役和修路成本；边镇得到的经验、损耗和未决情报会进入下一轮调度。朝廷可能将远征写成威慑或凯旋，地方承担者则更直接地面对马匹、车船、夫役和欠额。把战役只按胜败编排，会错过它怎样改变了河运、仓储与人事的优先次序。正是在这些不够壮观却持续发生的安排中，北京的北方位置变成国家财政与社会关系的一部分。

\subsection{材料把国家看成什么，也决定我们能看见谁}

洪武、永乐两朝的实录为这一时期留下密集的中枢年序：诏令何时颁发、官员怎样任免、军队何时出动、工程怎样被报告。它的密度不应使人误以为北京或南京看见了全部社会。实录中出现的云南、贵州、交趾、草原和港口，往往先以军令、贡物、战报或行政类别进入朝廷；妇女、佃农、军户家属、译者、船工与山地社区则较少留下能够与中央同样连续的记录。地方志、碑刻、契约、墓志和外部编年能把少数地点和关系带到近处，但它们又有保存、阶层和修纂的偏向。读者不应让一种材料代替另一种材料所没有看见的人。

云南的卫所、贵州的建置和交趾的直接统治尤其需要这种读法。中枢文书可以确认朝廷在何时派遣将领、设置机构或要求纳粮；一通桥碑、一次土司文书、某地后出的方志则可能保存工程、称谓或家族网络的一角。前者不能独自证明山地社会已经服从，后者也不能从一个地点推出整个西南。对地方行动者来说，新制度不只是名称变化：它可能改变道路的通行条件、调解诉讼的人、军粮的来源和可以寻求保护的对象。材料不足之处应当留下空白，而不是用王朝地图的颜色替代这些具体关系。

北方与海上亦然。朝廷将草原部众称为“虏”，将海上行动归入朝贡、禁令或“私出”的类别，首先说明国家如何处理威胁与准入，并不提供对方内部的完整自称。蒙古叙事、朝鲜与琉球的记录、航海者的见闻和港口遗址能提供不同的时间、地点与语言，却也不可以无缝拼接。船队抵港、使节抵京、互市或冲突发生，至少意味着某些道路和翻译关系曾经接通；至于这些接通惠及谁、伤害谁、持续多久，必须留给各类材料能够承担的尺度。

这种来源上的克制并不削弱明初故事，反而让它的压力线更清楚。朱元璋和朱棣的决定需要通过文书、仓储、驿路、河道和地方中介才会发生效力；同一条链在军队缺粮、河工受阻、山路难行或港口风向改变时也会显出断裂。国家能够更频繁地提出问题，并不等于它已拥有每一个地方的答案。把命令、报告、执行和后果分开，才能同时看见早明政权怎样扩展，又为何在每一次扩展中制造新的协调成本。

从应天建国到北京定都，明初并未完成一条直线式的“统一”。它把发令、登记、供给和战争分布在宫城、州县、港口、河道、山路和草原上；每一次扩张都使更多人和物资进入可调动的网络，也使道路、册籍、地方中介与资源征解成为下一次危机的条件。
\section{把建国送到道路尽头：1368--1387}

应天即位以后，最难的并不是再给新朝取一个名称，而是让一份在南京写成的命令能够穿过江北、山东、河北和西南的道路，成为仓、城、田与人的安排。北伐军接收大都时，需要接管的是城门、粮仓、官署和通往北方的驿路；旧元的官员、商人、工匠和居民并不会因改称北平而立刻拥有同样的处境。朝廷的诏书把归附写成一个可宣布的结果，地方接收则由谁保存账册、谁能催粮、谁被编入守城和运输来决定。史书较容易保存将领和城池的次序，对居民去留、市场停复和基层吏役怎样处理旧有关系，保留下来的痕迹远少得多。

北元向草原移动以后，\eventref{EM-1369-SHANGDU}{上都、应昌一线的追击}使这个差别尤其清楚。军队离开城墙进入水源分散、营地难以固定的空间，马匹、牧草、向导与补给比一张地图上的推进线更能限制行动。\eventref{EM-1372-LINGBEI-EXPEDITION}{三路北伐受挫}并不只是一次军事失利；它让南京必须承认，北方秩序不能由占领大都后的一次追击完成。卫所、边镇、互市、封贡和偶尔的远征以后会共同承担这种未封定的边界。明方文书能说明朝廷为哪一处粮道和兵马发令，蒙古材料与后出的汗统叙事则提醒我们，草原上的部众、首领和移动路线并不只是北京奏报中的背景。

海面构成另一条难以封定的边界。\eventref{EM-1371-HAIJIN}{海禁的颁行}把未经许可的船只、货物和来使归入需要约束的类别，却没有使东南港口停止面对季风、渔盐、借贷和区域贸易。对沿岸官员而言，禁令提供了何种船该问、何种人该查的语言；对船户和商人而言，航期、货源、船员和地方保护仍决定生计能否继续。遣使琉球、后来郑和的官船以及私人航行在同一片海上并存，并不相互抵消。朝贡文书最能说明谁被官方承认为使者，不能据此量出未入文书的贸易，或断言某座港口从此与外海隔绝。

\eventref{EM-1376-PROVINCIAL-COMMISSIONS}{三司分置}把新朝的谨慎写进地方机构。民政、军政和监察不再由原行中书省集中处理，看似把权力拆开，实际上也把协调任务送到更长的文书链上。征粮时，布政司要同府县、里甲和户部的口径相接；调兵时，都司、卫所和地方供给不能各自行动；发生诉讼和灾害，按察系统的查核又会影响哪一份报告先抵京。制度名称可以在会典和实录中查到，真正的行政速度却取决于书吏、驿传、船只与能够替各机构翻译地方情形的人。将三司写成一次性完成的“分权”，会遮蔽它日后不断要求协调的成本。

这种成本落到户与田上，才会显出国家能够看见什么。\eventref{EM-1381-LIJIA-HUANGCE}{黄册、里甲的编制}不是对安定社会所做的例行普查。元末战争造成的死亡、迁徙和逃亡已经打断许多家庭原有的居住与役属，国家希望把人户、军役和赋税重新接到可追查的格式。对于一户人家，登记可能意味着获得一个可据以申诉的身份，也可能意味着原先分散的差役被重新追到门前。书手、里长、业主、佃作者和军户家属怎样把婚姻、分家、欠债或迁居写进、写出或者留在册外，是中央表格难以完整保存的过程。

\eventref{EM-1387-FISHSCALE-REGISTERS}{鱼鳞图册等田土登记}进一步把界址、亩数与赋额连在一起。它为比对和催征提供工具，并不等于实际占有、耕作与收成已经被准确掌握。存世图册和徽州契约之所以珍贵，正因为它们偶尔让我们看见一块地如何在买卖、继承、抵押和税役之间被称量；它们也因保存地点集中、抄补层次复杂而不能代替全国。早明国家的扩展不是把乡村变成透明的数字，而是不断生产新的追索格式，同时也生产规避、协商与重新书写的空间。

西南把这些制度问题放进更长的行军线。\eventref{EM-1381-YUNNAN-CAMPAIGN}{明军由川黔进入云南}时，主将和朝廷可以制定会师、进军与设治的计划，山路、渡口、粮食与地方关系却决定队伍是否能继续前进。梁王、地方首领、城镇居民与山地社群不会以同一方式面对新军；抵抗、避让、供应和谈判都可能在不同地点同时发生。\eventref{EM-1382-YUNNAN-ANNEXATION}{昆明陷落后的卫所与建置}因而不是战争的自然尾声。军屯要有种子、驻军要有粮道、官署要有懂得地方语言和关系的中介，城市以外的道路和社区也要反复被纳入新的保护与征发关系。

云南的碑刻、方志和后来的土司材料能够使某些城、桥、寺观、军屯或家族露出痕迹；它们无法替所有没有留下文字的人说话。明方战报的“降附”说明朝廷希望记录怎样的结果，不足以证明各地接受了同一种治理。把这一点留在叙事中，不是削弱建国的规模，而是让读者看见开国军队与制度必须在真实的山路、河谷和地方生计中反复付费，才能把南京的命令变成持续的关系。

\section{继承战争怎样使用国家骨架：1398--1402}

洪武末年的继承不是一张已经写完的谱系图。太子朱标去世以后，皇太孙的名分与成年藩王掌握的封国、卫所和边地资源并存。\eventref{EM-1398-HONGWU-DEATH}{太祖去世、建文帝即位}把这层不稳定带到中枢：遗诏、宫中安排和诸王态度经过建文、永乐两次修史，不能由后来的胜负倒推为一致的阴谋或忠诚。较可确定的是，新皇帝需要让官僚、宗室、卫所和地方官相信谁能继续发令；而每一封削藩、调兵和征粮的文书都可能反过来改变他们的选择。

\eventref{EM-1399-JINGNAN-BEGIN}{燕王以“靖难”为名起兵}，使这个问题沿河北、山东、淮北和运河的道路展开。双方都以合法性语言说明自己，战争却首先考验谁能维持兵粮和信息。北平周边的卫所、渡口与仓储给燕军以起点；建文朝能够依靠的州县、河运和中央文书同样构成资源。白沟河、东昌和南下的反复并非只为产生一个胜败表：一条道路被军队反复使用，意味着民夫、船只、粮价、地方守备和治安都被卷进不易见于正史本纪的消耗。

\eventref{EM-1402-NANJING-FALL}{燕军进入南京、朱棣即位}之后，宫城火灾和建文帝去向成为后世最容易填入传说的空处。入城和改元可以确定，火中发生什么、建文帝是否出逃以及谁在何时作出决定，材料并不能许可完整复原。更重要的是，取得南京意味着取得印信、仓储、六部和重新整理前朝记忆的能力。永乐朝重修建文记录本身，提示我们胜者不仅赢得城门，也取得安排事件先后、定罪和沉默的制度位置。

内战留下的不是一个已经消失的“靖难问题”，而是一套更紧的北方军事与南方供给关系。谁能调动地方资源、谁能穿过机构之间的延误、谁必须向地方社会解释新的役令，都会在北京营建、北征和漕运重新显现。读者若只在南京城门前结束这一段，就会错过战争如何把洪武时期拆开的机构与账册再度拉成一条需要长期维持的国家骨架。

\section{北京的尺度与海外的尺度：1403--1424}

\eventref{EM-1403-BEIJING-NAMING}{北平改称北京}不是迁都已经完成的同义词。宫城、坛庙、仓储、城防、卫戍和粮道需要十余年营建；每一项工程都要在南方粮运、北方军需、工匠征发和地方财政之间排定先后。北京城建遗址能让人看到某些城垣、宫殿和坛制的物质阶段，不能替诏令说明是谁在何时作出决定，也不能计算每一名工役者承担的代价。南京保留留都地位，意味着双都并非旧都被抹去，而是文书、官员、仓储与仪式在更长距离上被重新安排。

\eventref{EM-1403-YONGLE-DADIAN-COMMISSION}{《永乐大典》的诏修}与北京工程属于同一时期的集中动员。抄写、校勘、检索、贮藏和呈进要求稳定的纸墨、人力与官署秩序；现存卷册能证明工程留下过的物质痕迹，不能把“收尽天下书”的宣称变成无遗漏或普遍阅读。国家希望决定哪些文字应被分类并留在中枢，地方书院、寺观、家庭和市场的读写生活仍由抄本、价格、师承和流通决定。知识汇集扩大了朝廷的文化可见性，也暴露出任何汇编都须依靠它未能完全控制的地方藏书与书写者。

\eventref{EM-1405-ZHENGHE-FIRST-VOYAGE}{郑和首航}让这种距离转向海上。船队从南京、刘家港出发时，诏书、赏物、翻译和仓储可以由官府组织；进入印度洋后，季风、港口准入、当地统治者、商人和翻译者决定一项使命如何继续。马欢、费信的见闻让港口的礼仪、物产和交易出现于文字，也带着随行者观看异域的角度。把它们同琉球、东南亚材料、港口遗址和碑刻并读，可以确认某些使节和航程，不能把船数、最大船型或贸易额写成一个无争议的总数。

\eventref{EM-1407-JIAOZHI-PROVINCE}{交趾设治}与海上使团形成尖锐的对照。前者试图将官署、卫所和税役直接延入红河三角洲与山区，后者更多依赖礼物、册封、港口和临时翻译来维持关系。两者都需要运输和中介，却把成本落在不同人身上。明方诏令能够说明设治的决定，越南编年与地方材料提醒读者，驻军的粮饷、地方抵抗和山地道路不会因名称改变而消失。后来的撤军并不使最初的设治成为无意义；它说明直接统治、海上交涉和边地经营各有必须反复支付的条件。

\eventref{EM-1410-FIRST-MONGOL-CAMPAIGN}{永乐亲征蒙古}、\eventref{EM-1411-GRAND-CANAL}{会通河疏浚}和\eventref{EM-1415-CANAL-TRANSPORT}{漕粮北运的稳定化}应放在一张连续的资源图中。草原远征需要车马、粮秣和消息，运河需要闸坝、河工、船只与沿线劳力；北京越接近北方战场，越依赖南方与河道送来的粮食。潘季驯以前的河工材料并不允许我们把早明河道写成永远有效的机器。它能说明朝廷努力维持某些河段，不能证明每一年都同样通畅，也不能把工费数字等同于沿河家庭付出的劳力。

\eventref{EM-1413-GUIZHOU-PROVINCE}{贵州设省}、\eventref{EM-1414-TULA}{土拉河交战}和\eventref{EM-1416-BEIJING-CONSTRUCTION}{北京持续营建}同时发生，使皇帝面对的不只是“扩张”或“防守”的抽象选择。山地的卫所与土司、草原的部众与道路、京师的木材与粮运都在请求同一套有限的文书和财力。\eventref{EM-1420-BEIJING-CAPITAL}{正式定北京为京师}后，\eventref{EM-1421-PALACE-FIRE}{三大殿火灾}把工程并未结束的事实送回日常。重建、北征、漕运和继承同时要求资源，\eventref{EM-1424-YONGLE-DEATH}{永乐帝北征归途去世}才使这副拉得很长的骨架交给下一代。死亡与继位可靠，秘不发丧和随行密议的细节多经后出叙事塑形；比传奇更能说明问题的是，洪熙、宣德必须在既有的双都、河运、边防和远征成本中重新排定优先次序。
\section{文书、案件与地方的沉默}

洪武朝在中枢收紧权力的方式，常被概括为废相和重典；若把目光放在文书实际经过的地方，便会发现集中裁断同时增加了信息的负担。\eventref{EM-1380-HU-WEIYONG}{中书省被废}后，六部直接面对皇帝，意味着谁汇总奏报、谁能说明某项财政或军政决定的后果、谁为延误负责，都更靠近宫门。近臣、书吏和部院并未因此失去作用，反而成为把千里之外的田赋、军情、河工和诉讼翻译成皇帝能够裁决的文本的关键。裁断权越集中，文本的筛选、摘要和传递就越可能改变一项地方事务被理解的方式。

锦衣卫、案狱和大诰将这种压力推进到官员与家庭。设卫、查案、颁布法令可以确认，不能从中推定每一处审判都同样发生。郭桓、蓝玉等案件所呈现的数字、同党和罪名，首先属于国家要惩戒、恐吓和重排行政关系的语言。对处理账册的官员和书手而言，严刑可能迫使他们更谨慎上报，也可能鼓励避责、迁延或把风险推给下层。对军功集团和地方精英而言，案件重排了谁能够接近中枢、谁可以成为可靠中介。官方文献保留了惩处的声量，较少保留被牵连家庭、部属与地方社会如何调整日常关系。

国家在法律上要求准确，在道路上却仍受距离限制。急递、驿传、渡口和仓储决定一份奏报到达北京需要多久，一道命令又要经过多少次转写才能抵达县衙。水旱、河患、盗匪和军队移动都可能使这条链变慢。早明的制度不能被看作已经完成的机器，而应看作不断需要人、马、船、纸张和地方保护来保持运转的基础设施。后世称为“中央集权”的成果，在当时往往表现为谁能承担传递成本、谁能解释延误、谁在缺少确切消息时仍要作出决定。

科举与学校也在这种基础设施上扩展。国子学、乡试、会试和荐举提供了把士人送进官僚体系的渠道，却不意味着乡村社会均匀拥有读书、赴试和获取文书的机会。南北榜案之后，籍贯、取士与地方声望进入宫廷讨论，显示选官制度不仅挑选个人，也重新分配哪些地方能通过文字进入中央。学校志、登科录和家谱能够留下部分成功者的路径，对未能读书、未被录取或因赋役和迁徙失去机会的人则较少发声。用一张榜单证明全国文化整合，会忽略知识与行政资源的地域差异。

宗教和营建提供了另一条可见的连接。寺观碑刻、桥梁题名、堤坝记事与地方祠祀往往记录捐资者、官员和工程名称，说明某种公共行动曾被组织；它们也通常强调功德、秩序与地方精英的自我呈现。洪武、永乐政府能以礼制和工程命令塑造国家的语言，不能直接规定一座村庄怎样理解祭祀、洪水或家族义务。正因材料各自有限，读者需要同时保留诏令、碑刻、契约与地方志的不同尺度，而不是让最完整的中枢记录替所有地方经验作答。

\section{双都、运河与家庭经济的长线}

北京成为常驻中枢后，南方并未只是向北输送一条无声的粮线。江南的粮食、布匹、木材、税银和工匠要经过县衙、仓场、河港、船户和闸坝才可能抵达北方；每一个环节都有损耗、借贷、季节和地方协商。运河沿线的家庭可能从运输、市场或工役中获得机会，也可能在河工、徭役和灾害中承担更大风险。官方漕运记录能够说明应运多少、从何处起运，不能自动说明何时实际到达，或谁替延误和损耗付钱。

双都格局也使官员、宗室和商人处于不同的空间选择中。南京保有礼仪、官署、仓储和文化资源，北京则成为北防和皇帝日常裁断的中心。两地之间的往来需要船、驿、文书和稳定的河道，任何一处堵塞都会让政治决定的时间变长。北京的工程、北征和宫廷礼仪要优先取得资源，南京及沿线地方也要处理本地的赋役、河患与社会秩序。所谓迁都的历史意义，不只是皇帝搬到北方，而是把整个帝国的财力和时间表重新校准到一条更长、更脆弱的供给链上。

永乐后期多次北征的结果尤其显示这种校准的限度。追击阿鲁台或瓦剌部众时，军队能否获得对方位置的消息，往往比皇帝的出发日期更决定行动。\eventref{EM-1422-SECOND-NORTHERN-CAMPAIGN}{第二次北征}、\eventref{EM-1423-THIRD-NORTHERN-CAMPAIGN}{第三次北征}所留下的不是简单的胜败，而是一次次将补给、道路、马匹和边镇承受力推到远处的试验。朝廷可以将出征写成威慑或恢复秩序，边地家庭和军户面对的则是持续征发、迁调和不确定归期。没有材料支持时，不应为他们虚构态度；但不能因为材料稀薄，便把他们从国家行动的成本中抹去。

从1368到1424，建国的连续性正来自这种看似琐细、却不断被重做的工作：给城门换上新的守备，给田亩换上新的登记，给河道换上新的闸坝，给一份诏令找到能送达的人。战争、继承、远航、设治和迁都都能在这张网络中留下转折；网络本身从未停止要求维护。下一代面对幼主、边警和地方冲突时，承接的并不是一个已经完成的帝国，而是一套能提出更多问题、也要求更多人承担回答成本的制度。
\section{远航所见的不是一片无主的海}

郑和船队出航时，朝廷以册封、贡使和皇帝威德来组织对外叙述；港口中的行动却要面对已有的城市、王权、商人和宗教网络。\eventref{EM-1405-ZHENGHE-FIRST-VOYAGE}{郑和首航}所能确认的是船队、使节和若干航路的展开，不能将途经地区写成等待明朝发现的空白。爪哇、满剌加、锡兰、古里等地各有自身的政治竞争与贸易规则，明使带去的礼物、军力和文书只是在这些关系中增加一种强大的变量。

船队的物质条件同样值得进入叙事。造船、征调水手、准备淡水粮食、辨认季风和修补船只，都依赖江南工坊、福建港口与沿海家庭的劳动。官方记录偏重朝贡物、异兽和外交仪式，较少记录工匠、翻译、医者和水手的报酬、疾病与死亡。不能以这种沉默虚构他们的个人遭遇，却应认识到“下西洋”并非一位使者或一组巨舰独自完成，而是陆海交通、官营采购和地方资源暂时被集中后的结果。

远航与海禁也不是简单矛盾。政府一面控制私人出海，一面派出规模巨大的官方舰队，表明关键不在于海是否可以通行，而在于谁能取得通行、交易和代表国家的权利。官方希望把海外交换纳入朝贡秩序，海上商人与沿海居民则会继续依照利润、季风和亲缘安排航线。后来的私贸和冲突不能直接归因于永乐政策，但国家许可与实际流动之间的落差，已经构成一个会长期反复出现的问题。

远航的回报也不能只按珍宝衡量。它让朝廷获得异域信息、仪式声望和对若干港口关系的影响，同时占用船材、粮食、兵员与行政注意。北方用兵、北京营建和运河修治并未因此暂停，财政与劳力必须在不同优先事项之间移动。永乐末年的国家是一种扩张中的秩序：它能把力量送到草原、江淮与印度洋，也因每一次远距离行动而更依赖那些容易被宏大叙事忽略的仓储、道路、河港和家庭。
\section{从军籍到仓单：北方经营如何穿过地方执行（1369--1424）}

洪武二年北进上都、应昌的队伍离开大都以后，北方的难题已不再是把一座城改名为北平。\eventref{EM-1369-SHANGDU}{上都、应昌一线的追击}逼着将领把行军看成一连串有先后之分的取舍：水源和牧草决定马匹能走多远，运来的粮食决定营地能停多久，而从前方折回的消息又决定下一批车马应当送往何处。皇帝可以在应天看到凯报，却不能从一封战报中取得草原上部众的完整位置。元顺帝北遁而死的消息，在《太祖实录》与朝鲜方面留下的线索中各有到达京师的路径；它能说明北元旧中心仍在移动，不能替远处的军队保证一条能把移动中的对手钉住的道路。\eventref{EM-1372-LINGBEI-EXPEDITION}{洪武五年的三路北伐}使这一限制不再只是战略上的谨慎。诸路未能按预想取得战果，并非一个可以由单一将领的得失解释的结局：从城镇仓储向草原展开的补给，要经过人、马、车和熟悉路线者，任何一环迟到或判断失准，原本互相配合的部署便会在广阔距离上散开。《明史》列传和后起的蒙古汗统叙事保存的责任判断并不相同；较可把握的是，夺取北方都城以后，朝廷面对的是一条需要常年经营的接触地带，而不是一条已经画完的边界。

这也改变了军籍和卫所应当怎样理解。把士卒编入卫所、把户口分入可供差遣的类别，确实让朝廷可以把边防需要转写成册上的责任；但一本在州县或卫所保管的登记，不能自行变成驻地的粮、马和可用人手。北方军户的家属仍要耕作、迁居、婚配和应付地方差役，远离原籍的调发则会使田地、役属与亲属关系重新错位。朝廷借军籍追问的是“谁该出役”，基层必须处理的却是“这个人此刻在哪里、由谁替补、粮从何处来”。《太祖实录》所记的调兵和设卫，使前一种问题留下清楚的日期；残存的户籍、契约和地方材料只零星保存后一种问题。不能因为零星就把军户写成一个没有差异的受命群体，也不能凭个别逃役或屯田记载推论全体军户的生活。它们至少说明，北方经营的成本不是在出征日才产生，而是在每一次核对人丁、交纳粮料和递送军情时落到地方关系之中。

\eventref{EM-1376-PROVINCIAL-COMMISSIONS}{三司分置}使这种成本有了新的行政走向。布政司要将赋役与民政文书向上汇集，都司要处理军政与卫所的调发，按察司则承担巡按、纠察和案件的另一条线；把职责分开，减少了单一机构同时握有军、民、财的可能，也使一项边防事务较少能够由一处衙门独力收束。假如某卫缺粮，地方官并非只要看见一纸军令便能解决：仓中有无可拨之粮、里甲能否承受新的催征、道路和船只是否可用、文书应报向何处，都会把布政、都司、府县和中央部院拉进同一件事。后来的《明会典》能够展示这些机构各自的名目与职掌，却不能把它们倒写成洪武九年以后每一省都立即形成同样顺畅的配合。正是在文书需要来回、责任必须彼此核对的时刻，分置机构才成为一种实际的行政经验：它令朝廷更容易追究哪一环失职，也令地方更难以用一条直线把命令变成结果。

洪武十四年的\eventref{EM-1381-LIJIA-HUANGCE}{黄册与里甲编制}，应放在这条军政和民政相接的链上阅读。战后迁徙并没有自然停止；一个家庭的原籍、现居、田地、军役和依附关系可能并不写在同一处。里长和书手向上递交的，是可以供征派和比较的户等，而不是一幅没有遗漏的乡村图景。对被登记者而言，入册既可能带来申明身份、厘清责任的凭据，也可能使早已离散或分家的负担重新被追索。黄册残卷与徽州一带契约让人看到，家庭往往同时以亲属、田产、债务和役属几种语言安排自身；这些文书不必与户部的栏目完全吻合。因而，朝廷在《太祖实录》中推进编制，能够确认国家要求各地完成何种登记，不能证明所有户都按同一日期、同一口径被纳入册中。国家能力在这里不是无所不见，而是能够反复要求地方把流动的人和资源译成它可以追问的格式。

六年后的\eventref{EM-1387-FISHSCALE-REGISTERS}{鱼鳞图册等田土登记}把另一种压力接了进来。军粮、赋税和地方行政需要可以落到土地上的责任，田界却会因买卖、继承、典押、荒废或迁居不断变化。图册把地块、相邻界址和赋额排列为能核对的记录，使县衙有可能把“哪一户应当负担什么”落到较细的空间；它却不能消除业主、佃作者、书手和里甲对同一块地的不同说法。鱼鳞图册在若干地区的保存，特别适合揭示登记怎样被书写、抄补和援引，而不适合被拼接为全帝国的实测土地表。把黄册与图册并读，可以看见早明行政并非先有一张准确的社会地图、再据此征发；恰恰相反，征发、诉讼、军役与地方协商迫使国家不断制作地图式的文本。北方边镇要求可调动的人，京师要求可转输的粮，地方册籍便成为连接两种要求的中介，也是矛盾最容易显露的地方。

永乐取得帝位后，北方的压力没有因宫廷易手而缩小，反而同海上行动在同一套调度中相遇。\eventref{EM-1405-ZHENGHE-FIRST-VOYAGE}{郑和首航}从南京、刘家港出发，首先意味着诏书、赏物、翻译、船员和给养须在陆地上被聚集、验收和装载；进入季风海域之后，它们又必须依赖港口的水源、修补、市场和当地许可。马欢、费信的记述之所以重要，正在于它们让使团看见的港口秩序、物产名称和翻译关系进入文字；它们并不是由北京向外延伸的航海日志，不能替各港口居民说明舰队来临如何改变生活，也不能用来确证船只规模与交易数量。官方舰队与海禁并置，显示朝廷关心的并非抽象的“通海”或“闭海”，而是谁有资格携带文书、贡物和武力通过海面。南京和刘家港的仓场把海上行动接在国家的采购与征调上，印度洋港口则提醒使团：海路的继续并不由出发地单独决定。

这一海上尺度使北方问题更值得重新衡量。\eventref{EM-1410-FIRST-MONGOL-CAMPAIGN}{永乐八年亲征}所面对的，不是等待京师挑选的固定敌军，而是本雅失里及相关部众在草原联盟中不断改变的位置。朝廷记下皇帝出行、会战和受降，蒙古方面的叙事则保留另一种对部众关系的排列；两边都不能把草原的移动压缩成一张静止的敌我图。皇帝亲征可以增强政治声望，也要求沿途仓储、车马和信息供给持续不断。若把同年的海外使团与北征只列为两项“盛举”，便会错过它们对于同一批行政资源的请求：一个要使礼物、文书和人员越过海洋，另一个要使粮秣、马匹和命令越过草原。国家并不是在两端各自拥有一套无限的力量；它必须靠已经登记的人户、可以征集的物料和能够抵达的道路，在相互竞争的优先次序中安排出发。

\eventref{EM-1411-GRAND-CANAL}{会通河的疏浚}于是不能只是北方政策的背景。山东河段的淤塞、闸坝与水势，决定江南粮食能否由船只接续北上；河工又要在农时、地方劳力、堤防风险和上级催办之间取得暂时平衡。北京所需的是按时抵达的米粮，沿线衙门面对的却是某一河段何时可通、哪一处闸坝须修、役夫和船户如何集合、仓中粮食怎样交割。工程命令与河渠志提供的是朝廷如何确认问题、怎样要求修治的连续线；运河沿线的题记和遗址调查能使特定工程地点变得可见。二者合在一起仍不能产生一张每年都准确的劳役和损耗账表。正因如此，河工最能说明行政执行的层次：它既不是皇帝意志的自动延伸，也不是“地方反应”四字可以概括的阻碍，而是许多人在水、土、粮、船与文书之间反复调整的过程。

到\eventref{EM-1415-CANAL-TRANSPORT}{永乐十三年会通河航路投入漕粮北运}，这条调整才以较稳定的转输形态同北京相接。所谓稳定，并不是运粮从此无须等待水位、维修或验收；它指朝廷得以把过去分散且高风险的北运安排，更多地嵌入可重复的河道、仓场和经手程序。对山东河段、江南起运地和京仓而言，同一石粮在不同地点意味着不同的责任：在起运处是征解和装船，在河港是交割与守候，在北京则成为军政和宫城的供给。官方的起运、报解和抵达名目不能互相替换，尤其不能直接当作沿线家庭实际负担的总和。可是，正因这些名目必须逐段对照，军户、里甲和地方官所面对的登记压力才会同运河连在一起。仓单、册页和闸口并非宏大政策之外的细节；它们是北方常驻中枢能够在粮食到达之前预先作出军事和营建决定的条件。

永乐十一年同年出现的两种远距离行动，更能显示行政链并没有单一的形状。\eventref{EM-1413-GUIZHOU-PROVINCE}{贵州设省}把卫所、驿道、州县和土司之间原已存在的关系置于新的省级名目下。黔中连接湖广、四川与云南的山路并不因为诏书而变得平坦；驻军的粮、文书的翻译、案件的裁断仍须由能够在具体地方行走和交涉的人承接。地方碑刻与《贵州通志》保存的建置线索，可以帮助辨认某些节点，却不能把不同山地社群的选择收束为“设省成功”或“土司服从”。同年\eventref{EM-1413-ZHENGHE-FOURTH-VOYAGE}{郑和第四次航行抵达忽鲁谟斯等港口}，所用的则是另一种中介：港口翻译、礼物交换、当地统治者与商人网络。两者都说明北京需要依靠自身无法完全控制的通道；差别在于贵州的官署尝试把责任留在地面上，海外使团则必须在对方已有的秩序中争取短暂的准入。把它们并置，不是说山路与海路可以互换，而是让人看见早明国家在不同空间内都要付出协调而非仅仅发令的代价。

\eventref{EM-1420-BEIJING-CAPITAL}{北京正式定为京师}后，这种协调有了更严格的时间表。南京仍为留都，官署、仓储和礼仪并未消失；北京却要求军政裁断、城防和常驻人口在更接近北方的地点获得日常供给。于是，一份来自南方的粮册不再只与地方赋役相关，也同京仓是否能支撑下一次调兵有关；一段运河的停滞不只让货物迟到，也会改变部院何时能够按已有储备作出决定。北京故宫的遗址和营建分期使这种中枢的物质尺度可以被看见，《明太宗实录》则保存了定都的朝廷程序；两类材料都不足以精确计算每一名工匠、船户或军户的成本。它们共同否定的，是把迁都想成一个命令落地后便完成的瞬间。定都真正把宫廷、边镇、仓场和地方册籍拉进了同一套持续校验的时序。

这种时序解释了\eventref{EM-1422-SECOND-NORTHERN-CAMPAIGN}{永乐二十年}和\eventref{EM-1423-THIRD-NORTHERN-CAMPAIGN}{随后北征}为何不能只以“寻得”或“未得”阿鲁台来评价。对手避开会战，首先使明军前方的消息、饮水和给养更加重要；对于京师和运河沿线，远征意味着已经被安排进仓储和道路的资源必须继续向北移动。明方行军报告能够确认出征和若干经过，不能把敌方位置、部众损失或威慑效果写成无争议的事实。太子监国与前线亲征之间的分工，至少表明远距离军事行动已要求中枢将日常裁断同皇帝在外的移动同时维持。\eventref{EM-1424-YONGLE-DEATH}{永乐帝归途中去世}使这一点突然变成继承问题：继位的事实可以确定，随行人员怎样隐瞒消息或各人内心如何盘算则不应由后出故事补成现场。更有解释力的是，洪熙政府接到的不是一份空白的权力，而是一整套已经在北征、漕运、双都、军籍和海外往来中运转的责任链。

从1369年的草原追击到1424年的北征归途，北方巩固并没有停留在战场边缘；它不断把军户的登记、田土的核对、仓粮的交割、河道的维修和文书的传递拉入同一条年序。海上的使团、黔中的设省与北京的定都并未使这条链走向同一种结果：有的关系依靠驻军和地方中介，有的依靠港口准入和翻译，有的则靠闸坝、仓场和经手程序维持。它们共同提出的限制是，朝廷能够扩大命令抵达的范围，却不能免除使命令在每一处真正执行的材料、时间与社会成本。早明国家的力量，正是在这些并不辉煌的核对、等待与接续中被制造出来，也在同样的地方显出它的边界。

还应把仓场看作这条链上最容易被名词掩盖的节点。粮食在县衙的簿册上成为应解之物，在起运港成为须由船户接收的货物，在会通河的闸口成为等待水势和验验的船载，到北京才成为可以拨给卫戍、工程或出征军队的储备。每一次转换都要求有人确认数量、封识和去向，也给误报、损耗、拖延与重新核查留下空间。河渠志和漕运条目保存的是朝廷如何把这些环节命名、督办和归责；沿线碑刻所保存的工程与捐资，则让少数地点的维修者和地方关系露出轮廓。二者不能合并为一份完整的实收账簿，但正因为相互不能替代，读者才不应把“南粮北运”理解成从江南自动流到京师的管道。\eventref{EM-1411-GRAND-CANAL}{疏浚}、\eventref{EM-1415-CANAL-TRANSPORT}{通航}和\eventref{EM-1420-BEIJING-CAPITAL}{定都}之间的连续性，恰好在这种逐段验收中形成：前一项工程为后一项转输创造条件，后一项转输又让北京能够把北方的风险纳入日常预算，而非只在战报抵达时仓促应对。

因此，早明行政的关键并不是某一种制度是否被一次颁布，而是不同记录能否在时间上接住彼此。军籍使卫所能提出对人的要求，黄册和田土登记使州县能把赋役与地方责任相连，仓单和河运文书又把这些责任转换成可向北移动的给养；但每一次转换都需要实际在场的书手、里甲、船户、工役、军官与地方中介。海上使团所带的诏书和礼物、贵州道路上转递的文书与军粮，也有相同的限制：文本可以确定一项权利或命令的起点，却不能保证它已跨过水路、山路和不同语言的社会关系。把这些环节置回1369至1424的年序，便能看见明初的北方巩固并非单靠边军或宫廷完成，而是由一次次登记、验收、转运和协商所维持；当其中任何一环失灵，远征、定都和对外行动的宏大目标都要重新回到地方执行的尺度上来衡量。

\section{仓场、递报与军籍：北方压力怎样落入地方日常（1368--1424）}

大都易手之后，北方并不是在一份改名诏书中才开始被经营，而是在城门以内外的物资去向中被重新组织。守城者和接收者首先要面对的，是旧有仓场是否还有可点验的粮、官署的印信和簿册由谁看守、城中与城外的道路能否让新派的军民往来。朝廷的军报可以把入城写成一个清楚的日期；它不替居民、铺户和旧吏说明次日的米价、雇役和生计。不过，正是这些未被整齐载入本纪的事务，决定北平能否从一次军事接收变成向北输送人马、向南递送文书的节点。仓不是单纯等待军令开启的容器。仓粮要有人清点、搬运、守护和交割，城门外的车马要有人验放，沿路的驿递又须有草料、夫役和可以暂歇的处所。一个地方若只在战报里以城名出现，未必因此没有承担这些工作；往往只是承担者没有以姓名进入中央的叙述。

北元向上都、应昌方向移动时，\eventref{EM-1369-SHANGDU}{北进的追击}把这种城市接收的后果向草原边缘推开。军队由有围墙、有库藏的空间进入水草与道路难以由一处官署预先安排的地带，前方所需的粮料不能只靠后方一次拨给。给养队、马匹、向导和返回的信使必须不断调整行进的节奏；一支队伍走得越远，后方越要决定哪些粮草应继续前送，哪些要留在城镇与卫所，以免前方未能合击时后路先空。关于元顺帝北遁和去世的消息，在明方实录与朝鲜方面留下的记录中并不以同一条线抵达我们面前。这种差异提示的并非某一边天然更可靠，而是边情本来就是经过使节、降人、商旅和军中报告才成为京师可以处置的消息。朝廷能够按接到的消息发令，却很难同时掌握每一支部众的去向、每一处水源的状况和每一段道路实际能承受多少车马。

\eventref{EM-1372-LINGBEI-EXPEDITION}{三路北伐的受挫}使这种信息与供给的间隔格外显形。后出的列传会归纳将领的功过，草原汗统叙事又会从另一种政治记忆排列对手；若把它们直接拼为一份毫无缝隙的行军图，反而遮蔽了当时最实际的问题。多路部署要求每一队在不同地点、不同天气和不同水草条件下作出相近的判断，然而前一日发出的计划抵达后方时，前方可能已因敌踪、牲畜或道路而改变位置。远征没有取得预想的战果，并不等于北方社会此后只剩防御；它使朝廷较少能把边界理解为一条由会战封定的线，更多须把城镇储备、卫所驻扎、互市机会、递报与偶尔的出征接成长期的应对。对北平附近的住户、军户和运输者而言，所谓“长期”尤其具体：不是一场战后便能取消的差事，而是何时要交草料、何时要修路、何时又有队伍经过的反复不定。

这种不定首先进入军籍。卫所名册把能够供差遣的家庭列作国家的兵源，却没有把一家人的劳力固定在原地。有人随军调往边地，留下的亲属要料理田土、赡养与本地役事；有人因婚姻、逃荒、投靠或屯种而离开原籍，册上仍可能保留一项尚待核实的责任。于是，军籍既是朝廷寻找人的办法，也是地方不断解释人为何不在、由谁暂代、田由谁耕作的场所。将军户简单写成专司作战的封闭群体，容易忽略他们同时生活在村社、市场和亲属网络中；把他们只写成受压迫者，又会抹去不同家庭借军职、屯田、地方关系或流动机会作出的不同行动。存世的户籍、契约和地方志只允许我们靠近少数地点的这些安排，不能推出一幅全国一致的军户图像；但它们足以提醒我们，兵员的可用性从来不是名册上一个永远不变的数字。

\eventref{EM-1376-PROVINCIAL-COMMISSIONS}{三司分置}之后，这些变动不会只停在卫所内部。都司收到调发的要求，往往需要府县、里甲或仓场配合粮料和道路；布政系统处理的赋役与民政，也会因军粮、迁徙和屯田发生新的交叉；按察官所受理的告讼与查核，则可能把基层的拖欠、冒名和争地送回另一条文书线。制度史中把三司列为并列机构很方便，实际的地方执行却常发生在它们职掌的交界处。某处军粮迟迟未到，是仓中没有余粮、运夫没有聚齐、册籍所载户丁已离散，还是不同衙门对责任各有说法？这些问题未必都会留下完整案卷，却决定一条看似明确的军令会在何处停顿。会典所呈现的是后来整理出的职责框架，不应倒过来作为每一省、每一卫在洪武年间已经无缝协作的证明。分置确实增加了相互制衡的可能，同时也让地方必须用更多往返的文书去确认谁该承担一项迟误。

\eventref{EM-1381-LIJIA-HUANGCE}{黄册与里甲的推进}使军籍的流动同一般户役的登记被放进可以相互比照的格式。朝廷需要知道谁可征粮、谁当出役、谁有田土；而家庭面对的却是分家、寄居、借贷、死亡、再婚和迁徙所留下的重叠关系。一户人家在乡里被认识的方式，可能同时包括族属、佃作、旧军役和近年的居住；向上报送时，书手却须将它们压入户等、姓名和责任的栏目。黄册残卷和区域契约所保存的，正是这种压入并不总能一劳永逸的痕迹。它们不能告诉我们每个北方军户怎样生活，也不应被拿来统计全国的逃役；它们让我们看到，国家的登记具有实践性，是在反复报送、核对、补写和争执中才取得效力。对地方官来说，准确的册籍能使催征有据；对被登记者来说，它也可能使原本因距离而模糊的义务重新变得可追索。

\eventref{EM-1387-FISHSCALE-REGISTERS}{田土登记的催成}又把这份追索落实到更细的地面。边地或近畿的军粮最终仍须从可以种植、征解和运输的土地中取得；一块田在鱼鳞图册上的界址、赋额和相邻关系，因而同远处的仓粮并非毫无关联。然而册上的地块不是一块永不移动的物：买卖、典押、继承、荒废和水患都会使使用者、业主和应役者不再完全重合。地方书手把田界写清，可能帮助县衙处理催征与争讼，也可能令某种此前尚可协商的负担有了更硬的依据。若只把图册当作国家视野突然完整的象征，就会错过它为何不断需要抄补、比对和援引。北方军事所需的粮，在这里经过的是田亩、里甲、县仓与河陆运输的一连串转化；每一次转化都给地方人留下解释、应付或重新协商的余地。

永乐初年北京被置于新的政治尺度后，北方的仓场问题变得更难从地方事务中分离出来。\eventref{EM-1403-BEIJING-NAMING}{北平改称北京}所启动的，不是把既有粮仓和道路自动搬到一个新名称之下，而是让宫城、卫戍、工地与前线更频繁地向同一方向索取供给。北京周边的驻军需要日常粮料，营建所需的木材、砖石和工匠也要等待不同路线上的人力和物资到齐；南京仍保有官署与仓储，意味着两地之间的调度并非单向替换，而是更长距离上的再分配。北京城建的遗存可以说明某些工程确曾展开，实录则使定名、调发和营建的朝廷次序可见；二者都不能精确折算每一次征调怎样落在某一户、某一船或某一座市镇。把这种未知保留下来，并不妨碍理解迁都的社会后果：越来越多本在地方循环的粮、材和劳力，要在新的时间表中为北方中枢留出位置。

\eventref{EM-1410-FIRST-MONGOL-CAMPAIGN}{永乐八年的亲征}使仓与关隘之间的联系进一步紧张。草原对手的部众、联盟和路线不会停在等待明军抵达的地方，前线将领尤其依赖能够逐级传递的侦报。一个消息若从边地传到行营，再由行营送回京师或沿路仓场，已经带着时间和叙述者位置的痕迹；朝廷据此作出的调粮、调马或改道决定，也可能在抵达时面对新的情况。明方记载中的受降、歼获和敌军人数，往往属于军功与朝廷威望的语言，不能替代草原社会内部的完整记录。较为确定的是，亲征并没有让北方获得一劳永逸的情报优势。它反而使递报、向导、翻译和边镇储备成为更显眼的资源：没有可靠的通道，兵力再多也难把遥远的行动者固定为可以决战的对象。

因此，\eventref{EM-1411-GRAND-CANAL}{会通河的疏浚}和\eventref{EM-1415-CANAL-TRANSPORT}{漕粮由此较稳定北运}应当从县仓与闸口的工作来理解，而不只从北京得到多少粮来理解。水位、淤积、闸坝、船只和验收程序使一石米在不同河段具有不同的风险。起运地的责任是把应解之粮集中并装载，河港要处理停泊、换船和守候，京仓则须按抵达的时日接收、保管和再行支给。三个环节在文书上都可以称为漕粮，却不能互相当作同一件已经完成的事。河工碑刻、闸坝遗址与实录中的工程命令能够使若干位置和政策节奏相互印证，但不能把沿线役夫的总数、每次损耗或地方家庭的负担化为精确的统一数字。更值得注意的是，运河把地方行政的时间感改掉了：一处堤防失修或水势异常，不只影响当地通行，还可能使京师、北征和营建的供给预期同时失准。

粮运进入京仓后也没有脱离地方社会。仓场需要负责盘验的人、搬运的人和维护秩序的人；粮价、临时雇役和船只停留会使城镇市场感到转输的波动。我们没有足够连续的材料去写每一座河港如何获利或受损，却不应把仓单当作粮食已经无摩擦地抵达皇帝手中的凭据。仓单是某一次交割被认定的证明，它的形成依赖丈量、封识、看守和责任分配，也可能引发催办、赔补和争讼。北京的常驻中枢之所以能够预作北方军事和宫城营建的安排，正在于它试图把这些散在不同地点的交割收进可预期的节奏；这种预期一旦遭遇河患、天气、盗扰或人手不足，又会把压力沿原路送回州县、村社和船户。

边情的递报与海上的联系在此并不属于两个互不相干的世界。洪武四年的\eventref{EM-1371-HAIJIN}{出海禁令}首先规定了沿海行动应当如何进入官府可识别的类别：谁携带文书，谁以朝贡名义往来，谁的船只因无许可而被视作需要查问。禁令本身不能说明海面因此沉寂；渔盐、季风、债务、亲缘和港口市场仍会让人以各种方式面对外海。朝鲜、琉球方面的文书之所以重要，不是因为它们能替明廷补写所有沿海社会，而是因为同一类来使、馈赠和消息在不同政权的记录里常有不同的到达路径。它们让我们看见海上信息如何穿过语言、历法和文书格式，也提醒我们不可把朝贡的名称、私人贸易的实际和沿海防务的顾虑混为一个无差别的过程。

到\eventref{EM-1405-ZHENGHE-FIRST-VOYAGE}{郑和首航}时，官方舰队的起航把这种“谁能通行”的问题推到更远处，但并没有取消港口和中介的限制。南京与刘家港能够集结诏书、赏物、船员和粮水，不能替船队决定每一个港口是否准入、何种翻译可用、何种交换能被当地接受。马欢、费信留下的见闻让一些港口的物产、礼仪和地名进入明代文字，所呈现的始终是随行者经由使团框架得到的视野。对早明地方社会而言，更直接的后果常在出海之前：造船所需的木料、帆索和工匠从哪里取得，船员与后勤人员怎样被征集，沿海官员怎样在禁令与官方任务之间重新分配注意力。我们不能因材料缺口替这些人安排感受，却可以看到，远航的可能性来自陆地上的仓场、工坊、港口和文书先被接合起来。

北京定都以后，边情、漕运和海上消息更清楚地显出同属一个有限的行政视野。\eventref{EM-1420-BEIJING-CAPITAL}{北京成为京师}，并不使南京、山东河段或东南港口只剩被动供给的角色；相反，它们必须不断以粮、物、人和消息参与北方中枢的日常判断。\eventref{EM-1422-SECOND-NORTHERN-CAMPAIGN}{再次北征}与\eventref{EM-1423-THIRD-NORTHERN-CAMPAIGN}{次年搜寻阿鲁台部众}没有立刻化为决定性会战，恰能显示信息不是军事行动的附属品。敌方避战时，前线最缺的未必只是兵力，而是关于水源、路线、部众位置和补给可能性的及时判断；后方最难的也不只是继续发粮，而是判断应把有限的粮、马和车辆押往何处。实录能够按朝廷所收的报告排出出征与回师的顺序，不能令我们看见每一条未经上报或未被保存的边地消息。把不确定性写入叙事，才不会把草原行动误写为地图上两支固定军队的相遇。

\eventref{EM-1424-YONGLE-DEATH}{永乐帝北征归途中去世}时，这些联系没有随一位皇帝停止。死亡和继位是可以钉住的节点，关于随行人员如何保密、谁曾说过什么的戏剧化细节则多出自后来的叙事，不必填入一个仿佛可见的现场。真正留给洪熙朝的是一组已经相互牵动的事务：北方的卫所和边情要继续有人递报，京仓和运河要继续有人核验，军籍与里甲仍要面对迁徙和欠役，沿海的官员也仍要在控制通行与处理实际联系之间作出判断。从1368到1424，国家的扩张并非只表现为更多的城名、官名和战役；它还表现为越来越多的家庭、仓场、闸口、驿站和港口被要求把各自的时间接入朝廷的时间。这个过程使北京能够向更远处作出选择，也使每一次选择都更依赖地方能否承受、解释并接续它的成本。

\section{县门、田界与案卷：早明的秩序怎样进入乡里}

若只从南京的诏令看早明的法与政，县衙往往像是一个没有厚度的终点：命令自上而下抵达，田粮、户役和案件再自下而上汇成数字。实际的县门却是两种时间相互顶住的地方。上级要的是可以按期核对的名目，来诉讼、应役或请求更正的人带来的却常是已经变动的家计、田界和亲属关系。一纸牒文到了县中，未必立刻化成判决；它可能先由书吏抄录、由差役传唤、由里中辨认，又在当事人补呈的契约、保结或相邻人的说法之间停顿。反过来，一场田界或役属的争执，即使始于村旁的小路、田埂或一户人的迁居，也可能因为需要查验旧册而被送入国家能够追问的文字。早明官府所建立的，并非一条永远笔直的裁断线，而是一套使地方关系反复变成可问、可查、也可被重新叙述的路径。

\eventref{EM-1376-PROVINCIAL-COMMISSIONS}{洪武九年分置三司}，为这条路径规定了不再完全相同的出口。布政司接续民政、赋役与地方行政，按察司承担纠察和司法监察，都司则处理军政；一项涉及人户、田土与军役的纠纷因而不必天然只属于一处衙门。这样的分置首先是一项防范集权于一人的安排，却也把地方官置于更复杂的责任网络中。县官面对一桩拖欠，须分辨它是因户主逃亡、田地易手、役法不明，还是征派本身超过乡里此刻所能承担；若牵涉军户、屯田或军粮，文书又可能越过民政常轨而与卫所相接。按察系统能够追问失察和枉断，不能替府县取得最初的事实；都司能够要求人力和粮料，不能单凭军令弄清某一户为何已迁离旧籍。制度将权责拆开，未将事实自动拆清。后来会典所呈现的职掌边界，是理解这些机构的必要入口，却不能证明洪武九年之后每个府县都已按同一节奏把案件转送、覆核和结案。

这种差别首先关系到谁能在县门前被认作可以说话的人。州县官未必亲见每一块田、每一次分家或每一户的出入；他所见的是由书手、里长、邻人和当事人共同组织出来的说法。里长的地位并不使其天然代表全里人的利益，他既要协助催办，也会置身于熟人关系、债务往来和地方声望之中。书手能够把口头陈述排进可上报的格式，却也受限于所持旧册、官府时限和上级对责任的追索。相邻者关于田界的证言，可能照亮一段共同使用的水沟、坟地或田埂，也可能带着自身的租佃、买卖和亲属牵连。早明法律与行政因此不能只被理解为皇帝和官员对“民”的单向规训；它还迫使地方社会不断决定谁可作保、谁的记忆可被写下、谁应为一项未能说明的义务承担后果。材料很少允许我们重建某个无名当事人的完整经历，但正因如此，更不能把没有留下姓名的大多数人写成没有判断和行动的背景。

\eventref{EM-1380-HU-WEIYONG}{胡惟庸案后废中书省}，又改变了地方信息如何被想象为中央可以裁断的事务。六部直接对皇帝负责，并不意味着一件来自县中的案情就能原样进入宫廷；相反，层层摘要、转呈与核覆变得更加关键。地方官若担心某项解释日后被视作隐瞒，便有更强的理由要求册据、具结和多方证佐；下层承办者也可能将模糊的实际处境整理成较易向上交代的门类。这样的压力不必然带来较准确的判断。有时它会迫使旧账被重新比对，有时也会让责任沿着衙门内部向下推移，使最接近田地和户口的人承担无法证明的缺漏。案狱叙事尤其提示这种危险。《太祖实录》能够定位制度变动，《大诰》和后出的传记则保存朝廷如何界定罪名与警戒；其中关于结党、交通的指控不能脱离形成它们的审讯与政治语境，被当作已经独立证实的地方事实。把这一限度留下，才能同时看到严厉惩治如何重排官僚关系，又不把案卷的语言误当作所有相关者的内心和行动。

到了\eventref{EM-1381-LIJIA-HUANGCE}{黄册、里甲编制推进}的时候，所谓“查明”获得了一种格外具体的物质形式。纸张上的户名、户等、田土和役属，使官府能够把本来分散在村落、亲属与田界中的责任置于同一套比照中。可是家庭并不会为了便于登记而停止变动：有人因饥荒或婚姻寄居，有人分家后仍共耕旧田，有人名在原籍而劳力已在别处，有人持有地业却将耕作交给他人。面对这些情形，造册不是把既成事实抄写下来，而是要求地方把不易并列的关系暂时归入一个能供催征和核验的名目。对有田产或稳定依附关系的人，入册可能提供主张权利、说明役属的凭据；对迁徙者、寄居者和劳力不足的家庭，同一册页也可能使旧有责任在新的地点重新显形。国家获得的不是完整无缺的乡村肖像，而是一种可以反复要求乡里回答“此人、此地、此役现在如何相连”的方法。

现存黄册残卷之所以比整齐的制度叙述更值得细读，正在于它们留下了这种方法的边缘。徽州一带的册页与契约常使人看见户名之外的田产转手、亲属分合和债务关系；它们可以让研究者追索少数地方的书写习惯，却不能充作所有府县的缩影。存下来的文书本身有选择：易于保存的往往是能进入家族、业主或官府保管的纸张，而不是一次没有写成字的争论。契约也通常服务于买卖、典押或分产的特定目的，不会替我们记录每一位佃作者、短工或寄居者怎样理解一项赋役。正因证据的可见范围如此有限，较稳妥的读法不是把黄册当成精确的人口表，而是把它当成一次制度要求与地方关系相遇的痕迹。它显示国家希望怎样排列人户，亦显示地方人在何种时候必须以田地、亲属和旧约来回应这一排列。

这种回应并非都以公开对抗的形式出现。一个家庭可以设法补出可以承担役事的人，也可以借助亲属、业主或里中关系解释暂时的缺额；地方中介可以催促、担保或拖延，县衙则可能在需要尽快完成上报与避免将来追责之间取舍。我们很难据此为早明乡民概括出统一的“顺从”或“抵制”，更不能把后世关于逃役的概念不加区别地投回每一份早期册籍。能够把握的，是登记使原来可在地方关系中缓冲的义务更容易被比较，也使失去劳力、田地或居所的家庭更容易被问到旧有责任。它因此同时增加了官府的追索能力与乡里的解释负担。某些义务因被写入册页而更难完全隐没，另一些义务又因册页与现实脱节而催生新的报称、纠纷和复核；两种结果不能预先由法令文本决定。

\eventref{EM-1387-FISHSCALE-REGISTERS}{鱼鳞图册等田土登记的催成}，把这种追问从人户进一步压到相连的地面。田界不是地图上的抽象线，而可能是一条沟渠、一处坟地、一棵标识边界的树，或两家长期共知却未必每年重说的分界。县衙将地块、界址和赋额编入图册，意在使赋役有可核的落点；一旦地被买卖、典押、继承、荒废或受水患影响，原有的落点又会遇到新的占有与耕作关系。谁在册上有名，谁实际耕种，谁最终能够交纳，未必始终是同一个人。图册能够为田争、催征和责任划分提供一份有分量的参照，却不能消灭围绕参照本身发生的争论。地方书手写下界址，并不是在空白土地上创造事实；他要选择怎样容纳已有的称谓、相邻关系和旧约。业主、佃作者与相邻人也可能从同一块地看到不同的权利与风险。

鱼鳞图册和相关契约因而让“司法”显出比刑名更宽的含义。许多影响一家生计的案件未必以惊人的罪名出现，而是围绕谁该负担、谁可证明、哪一份旧文书仍有效展开。县官必须在催科的急迫、上级的追问和地方关系的持续性之间作出判断；若断得过快，可能将某种暂时的困境固定为长期负担，若迟延不决，又要面对积欠、上诉或失察的风险。按察官和上级衙门可以为复核提供另一条路线，却也会增加来往文书、候讯和查验的时间。没有连续案卷的地方，不能凭制度文本替任何一方写出胜诉或冤屈；然而徽州的图册、契约及其抄补痕迹足以表明，田土并不是国家单方面宣布赋额即可稳定控制的资源。它必须不断经过地方的记忆、书写和争执，才能被保持为一个可征、可断、可再核的对象。

从三司分置到黄册、图册的编成，早明官僚体系所获得的，不只是更多名目，而是把互不相同的地方问题接进较长程序的能力。县中一项役属的疑问可以牵出里甲的说明，田界的变化可以牵出旧约与相邻证言，军役或赋粮的缺额又可能使府县、卫所与省级机构彼此问责。程序有时给弱小或远离权势的人留下申明和复核的缝隙，有时也会使他们更难摆脱已经写定的身份；两种可能都不能由“中央集权”一词替代。中央的确能通过文书把远处的问题变成可裁量的事项，但每一次裁量都依赖有人愿意或不得不提供材料，也依赖县门之外的社会关系能够承受不断的查验。早明秩序的力量正在于它把这种查验做成可重复的日常；它的界限也正在于，纸上的一致性永远要同迁徙、分家、田界、债务和地方记忆重新相遇。

在这一过程中，地方官并非只是把朝廷的分类照抄给乡里。他必须在有限任期内同时处理钱粮、治安、诉讼和催办，所能依靠的是一批熟悉本地人地关系、却未必完全受他控制的承办者。县衙需要书手认读旧册，需要差役找到被传唤的人，需要里中有人说明一段关系何以改变；这些人提供的便利，也构成官府最难完全审查的环节。上级要求限期造册或清理积欠时，县官有理由把不合格式的事实催成整齐的说法；乡里则可能试图让自己的说法保留更多余地。两方都不是抽象制度的被动部分。前者要在失察与苛催之间保存仕途和政绩，后者要在供役、家计和邻里评价之间维持可过的日常。材料多保留呈送到官府的一面，较少保留没有走到纸面上的协商；这并不使后者不存在，只要求我们不以一份完成的文书取代它。

黄册所带来的一个深刻后果，是让时间本身成为可以被追究的对象。某人何时离籍、某田何时转手、某项役事从何年开始拖欠，都不再只是当事人各自记忆的事，而会在新旧记录的对照中变成官府需要判断的先后。可是，地方生活并不总以官府要求的年度节奏变化。灾害、疾病、婚姻和收成可以在一季之间改变一个家庭的劳力，旧册却可能仍按上一轮编造时的关系提出责任。于是复核并不只是技术性的更正：它关系到一户人家能否证明自己已无力承担，关系到同里的人是否愿意承认其离散，也关系到经手者愿不愿意将不便上报的差额写出来。《太祖实录》所见的是朝廷推进编制的年序；黄册残卷让人看到的是编制之后仍会留下空格、增改和不齐的现实。两种材料的距离，恰是理解早明治理不能省略的部分。

田土也是如此。一片地的价值不止是赋额，还包含灌溉、通行、祖产、租佃和邻接关系；其中任何一项变化，都可能让册上原本清楚的归属在实际使用中变得可争。鱼鳞图册把边界编进连续的地块次序，给县衙提供了比口头记忆更可复查的凭据，却不必然让每一条边界更少争执。有了可以援引的图册，当事人或许更能指明自己主张的依据；同一图册也可能让不熟悉书写、无力保存旧约的人处于不利位置。我们没有足够材料衡量这种不利在各地如何分布，不能把徽州所见的文书习惯推广到全国。不过，正因为图册与契约在少数地区能够相互照见，才可以确认国家的田土核算必须同地方已有的产权语言打交道，而不是从一张空白纸上重新发明所有权利。

司法的社会后果还在于，它改变了人们求助于谁。过去可以由亲属、邻里或地方权威暂作调停的冲突，一旦牵连役属、田额或官府要求的保结，便可能需要进入县衙的程序；进入程序并非必然得到更公平的结果，却会使争执的用语、证人和期限受到新的约束。相反，许多本应能够上达的冤抑或困境也可能在路费、等待、文书能力和地方势力面前停住。不能把没有案卷保存理解为没有冲突，也不能凭制度中允许申诉的条文断言每个人都能抵达裁断者。较有把握的是，三司分置和中枢对责任的严密追问，使地方案件更容易被放进逐级核覆的想象中；这一想象本身已改变官员、书手和乡里中介处理风险的方式。它要求每个人预先考虑：若日后被问起，自己手中能拿出什么文字、什么人证、什么旧例。

因此，早明的登记与审断不应被分作一件关于财政的事和一件关于刑名的事。户名与田界一旦成为可供比对的依据，便会进入催征、争产、役属与失察的不同问题；一次判处或一次追责，也可能反过来促使乡里在下一次造册时改变申报的方式。这里并不存在一条从法令到服从的单向因果，而是一种持续的反馈：官府愈需要可核的资料，地方愈须把变化说成可被接受的形式；地方愈以新的情形冲击旧有格式，官府愈可能要求再查、再造、再断。洪武时期的制度建设由此留下的不只是高悬的法条和整齐的机构名称，更是一种深入地方日常的压力。它能使某些责任不再轻易被遗忘，也会把原本可以缓冲的生计危机转成必须在文书中说明的欠额与过失。对这一压力的承受并不平均，而现存材料也不足以替那些沉默的人排出一张完整的受益与受损名单；保留这份不完整，正是避免将早明地方社会写成制度图例的必要条件。
